\documentclass[11pt]{article}
\usepackage[utf8]{inputenc}
\usepackage[T1]{fontenc}
\usepackage{lmodern}
\usepackage[margin=1in]{geometry}
\usepackage{csquotes}
\usepackage[style=authoryear,backend=biber,sorting=nyt,maxbibnames=99,giveninits=true,uniquename=init]{biblatex}
\usepackage[hidelinks]{hyperref}
\addbibresource{references.bib}
\begin{document}
\section*{Nitrate Export: Magnitudes and Biophysical Controls}

Nitrate is the overwhelmingly dominant nitrogen species leaving artificially
drained cropland, and its export is one of the best-quantified consequences of
subsurface drainage. Across continuously monitored plots, drainage districts,
and watershed outlets, three findings recur with strong support: the
subsurface pathway carries the great majority of nitrogen loss, the magnitude
of that loss is governed first by hydrology, and the underlying
concentration--discharge behaviour reflects rapid, macropore-mediated delivery
from largely organic-derived nitrate sources that are only partly attenuated by
denitrification. This section establishes those magnitudes and biophysical
controls---the ``what and why'' of nitrate export---and couples load to flow
specifically for nitrate; broader flow-regime behaviour is treated in the
hydrology section, and all management levers are deferred to the management and
design sections.

\subsection{Subsurface drainage as the dominant nitrogen pathway}

\subsubsection{Pathway dominance}

Multiple independent studies establish that in tile-drained landscapes the
subsurface pathway carries most of the nitrogen. In intensively monitored
Norwegian catchments, subsurface drainage carried $76$--$91\%$ of runoff and
$87$--$95\%$ of total-N (TN) loss, with drainage TN concentrations $2$--$4$
times those in surface runoff and nitrate-N constituting $71$--$85\%$ of TN
\parencite{Bechmann2014}. New Zealand mole-pipe systems captured about $62\%$
of surplus rainfall as drainage and routed roughly $80\%$ of soluble N, while
surface runoff dominated P, sediment, and \emph{E.\ coli}
\parencite{Monaghan2016}. At watershed scale, tile drainage accounted for
$56\%$ of annual discharge and $44$--$82\%$ (mean $62\%$) of NO$_3$-N export in
Ohio \parencite{Williams2015}, averaged $57\%$ of streamflow and dominated
nitrate export in Indiana \parencite{Grose2022}, and in Iowa the largest drains
contributed $70$--$87\%$ of exported water and nitrate \parencite{Coupe2020}.
In flat Illinois terrain, subsurface stations lost
$13$--$30$~kg~ha$^{-1}$~yr$^{-1}$ of nitrate-N ($16$--$26\%$ of applied N), with
more than $70\%$ of samples exceeding the $10$~mg~L$^{-1}$ maximum contaminant
level (MCL), while surface runoff was negligible \parencite{Dattamudi2020}.

\subsubsection{Loads, concentrations, and MCL exceedance}

Reported magnitudes span a wide but coherent range. Field- and watershed-scale
nitrate-N yields commonly reach or exceed $20$--$60$~kg~ha$^{-1}$~yr$^{-1}$ in
maize--soybean systems
\parencite{Gentry2009,Gentry2014,Dattamudi2020,Tomer2008}, and tile
concentrations frequently approach or exceed the $10$~mg~L$^{-1}$
drinking-water limit. In three Iowa drainage districts tile water averaged
about $13$~mg~L$^{-1}$ NO$_3$-N and exceeded the standard $85$--$90\%$ of the
time \parencite{Schilling2012b}, and across three catchments the flow-weighted
mean was $17.1$~mg~L$^{-1}$ with $74.2\%$ of samples above the MCL
\parencite{Ikenberry2014}. New Zealand tile drains exported
$14$--$109$~kg~ha$^{-1}$~yr$^{-1}$ TN ($58$--$90\%$ nitrate)
\parencite{Tanner2011}; Illinois riverine and tile yields averaged $30$--$32$
and ranged to $57$--$59$~kg~ha$^{-1}$~yr$^{-1}$ \parencite{Gentry2014}; and very
high concentrations characterise some European and irrigated systems, with
Italian drainage averaging $94$--$216$~mg~NO$_3$~L$^{-1}$ \parencite{ROSSI1991}
and Polish water usually exceeding $70$~mg~L$^{-1}$ \parencite{Dragon2016}. A
key methodological finding is that interannual variability in load greatly
exceeds that of concentration: across 29 Iowa watersheds, flow-weighted
concentration ($0.6$--$15.5$~mg~L$^{-1}$) was a more robust indicator of
management change than load or yield \parencite{Danalatos2022}. Early Iowa work
already characterised nitrate loss as ``highly variable and strongly dependent
on the amount of water lost'' \parencite{BAKER1975}, and long-term records
document real declines: German drainage N discharge fell from $33$ to
$10$~kg~N~ha$^{-1}$~yr$^{-1}$ between the 1980s and 1990s, driven mainly by a
reduced N surplus \parencite{Hirt2005b}.

\subsubsection{Fertilizer input, hydrology, and soil-N legacy}

At regional scale, fertilizer input and watershed hydrology jointly regulate
export. About $34\%$ of applied fertilizer N was exported to streams in the
Mississippi basin, wetter watersheds exported a larger fraction, and tile
drainage acted as a hydrologic modifier that increases throughput and shortens
residence time, raising nitrate loss to the Gulf of Mexico
\parencite{Raymond2012}. Direct comparison confirmed the causal link: nitrate
load from an unfertilized Minnesota field was $58\%$ of that from a fertilized
field \parencite{Ghane2016b}. Multivariate models combining water yield, prior
corn yield, and fertilizer sales explained $96\%$ of annual nitrate-yield
variation \parencite{Gentry2014}. Losses stem partly from asynchrony between
mineralization and crop uptake and from large mineralizable soil-N pools, so
land-management effects may only become visible after years to decades owing to
soil-N legacies \parencite{Cambardella1999,Bauwe2020}.

\subsection{Hydrologic control of nitrate export and seasonality}

The single most robust cross-scale finding is that nitrate load is controlled
by drainage discharge. Annual nitrate-N loss relates strongly to annual flow,
with differences among drain spacings explained by flow
\parencite{Kladivko2021,Bauwe2020}; in Italy NO$_3$ loss correlated with
drainage water ($r=0.93$) but only weakly with precipitation ($r=0.45$)
\parencite{ROSSI1991}. Because export tracks flow, interannual variability is
driven principally by precipitation and drainage volume rather than by
concentration \parencite{Danalatos2022,BAKER1975}.

High flows dominate export. More than half of Iowa catchment nitrate export
occurred during the highest $10\%$ of daily flows \parencite{Ikenberry2014},
and $79$--$95\%$ of Indiana NO$_3$-N export occurred at the 60th--100th flow
percentile \parencite{Speir2021}. Because loads track flow they are strongly
seasonal, concentrating in the non-growing season when evapotranspiration is
low and drainage is high: about $70\%$ of annual drainage and nitrate load
occurred in the November--April fallow season in Indiana
\parencite{Kladivko2021,Kladivko2004}, and $81\%$ of Ohio tile discharge
occurred in the nongrowing season, so nitrate and TN loads were greater then
despite higher growing-season concentrations under corn \parencite{King2016}.
Short episodes can dominate: $21$~kg~N~ha$^{-1}$ was lost in only $11$ days
during snowmelt in Japan \parencite{Hayashi1999}. The broader flow-regime
behaviour underlying this coupling is treated in the hydrology section.

\subsection{Concentration--discharge behaviour, sources, and denitrification}

\subsubsection{Concentration--discharge behaviour and rapid response}

High-frequency records reveal characteristic concentration--discharge (C--Q)
signatures. Annual nitrate export is largely chemostatic---yields rise
roughly proportionately with runoff---and persistent chemostasis has been
documented in Illinois and Ohio tile drainage
\parencite{Speir2021,Hu2024,Williams2015}. Storm-event hysteresis is
predominantly counterclockwise ($77\%$ of Indiana storms; $>75\%$ of a
long-term Ohio record), indicating dominance of distal nitrate sources during
delivery \parencite{Speir2021,Liu2022b}. Analysis concluded that intensive
subsurface drainage, rather than deep groundwater, was the dominant factor
shaping nitrate C--Q relationships \parencite{Liu2022b}. Tile nitrate responds
within hours of water inputs: after irrigation it rose sharply within hours
then dropped abruptly, with preferential flow intercepted by the drain
originating close to and above the tile \parencite{Mohanty1998}, and bypass
(macropore) flow carried most leached N in Japan \parencite{Hayashi1999}.
Threshold behaviour is common: tile discharge began once the soil-moisture
deficit exceeded $-17\pm3$~mm and the water table rose above $0.84\pm0.03$~m
\parencite{Williams2022}, and antecedent moisture modulated discharge and
solute pulses \parencite{Trentman2020,Kennedy2012}.

\subsubsection{Sources, transformation, and denitrification}

Dual-isotope studies indicate that drainage nitrate is often derived largely
from organic sources (manure and soil organic matter) via nitrification, with
denitrification not always dominant. In Ontario, $\delta^{15}$N and
$\delta^{18}$O relationships lacked the 2:1 denitrification slope, pointing to
organic sources and continued nitrification through all seasons
\parencite{Smith2011,Smith2011b}, and rapid nitrification of applied ammonium
($>90\%$ within $72$~h) supplies ready nitrate for leaching \parencite{Hu2024}.
Where denitrification is active it is detectable: along Polish flow paths
nitrate declined while sulfate, excess N$_2$, $\delta^{15}$N, and
$\delta^{18}$O increased \parencite{Dragon2016}, and nested French watersheds
implied $34$--$82\%$ nitrate retention dominated by underground denitrification
\parencite{Billy2011,Billy2013}. In Iowa drainageways, denitrification in
slowly permeable soils could remove essentially all upland groundwater nitrate,
functioning as a natural treatment system
\parencite{Schilling2007,Cambardella1999}. Yet where drains intercept water
before it reaches these attenuating zones, export is high: tile discharge is
dominated by pre-event water---three-component separation gave $65\%$
groundwater and $29\%$ shallow soil water---yet tile drains carried the highest
NO$_3$-N of any source and contributed over $99\%$ of watershed nitrate export
\parencite{Williams2022,Grose2022}.

\printbibliography
\end{document}
